\scanpage{193}\mspage{49}
% archive sheet number: 193

$u(\unsure{r|t'})(\gamma\unsure{\alpha'|\alpha})=u(r)(\gamma\beta)$, donc $u(\unsure{r})\gamma$ égalise $(\alpha,\beta)$, qed

La vérification de d) \unsure{en|un} \unsure{AQT}.

4. \uline{Foncteurs biasphériques}

Soit encore

\[
u:I'\longrightarrow I
\]

une flèche de $\mathrm{Cat}$, soit $M$ une catégorie, d'où \deleted{où les lim} \deleted{supposons que}

\[
u^*:\ M(I)\longrightarrow M(I')\ ;\tag*{(4.1.)}
\]

\deleted{\unsure{pour qu'}} nous cherchons des conditions \inserted{sur $u,M$} pour que ce foncteur soit pleinement fidèle. Si le foncteur adjoint à gauche $u_!$ (resp. adjoint à droite $u_*$) existe, \unsure{On|on} sait qu'il revient au $=$ de dire que le morphisme d'adjonction

\[
u_!u^*F\longrightarrow F\tag*{(4.2)}
\]

resp

\[
u_*u^*F\longleftarrow F\tag*{(4.\unsure{2|1} bis)}
\]

soit iso pour tout $F$ dans $M(I)$. Disons alors que $u\deleted{\unsure{f}}$ est $M$-connexe, ou $M$-biasphérique, cette terminologie sous-entendant que c'est là une propriété qui se réfère au dérivateur $\mathbb{D}_M=S\longrightarrow M(S)$ sur $\mathrm{Cat}$ défini par $M$. Notons que

\[
\begin{aligned}
&\text{Si }u\text{ est }M\text{-}\deleted{\unsure{\mathrm{bi}}}\text{connexe, alors il est aussi}\\
&\deleted{\unsure{Hom}} M(S)\text{-connexe, \unsure{quel|ceci} pour }\deleted{\text{tout}} S\text{ dans }\mathrm{Cat},
\end{aligned}\tag*{(4.3.)}
\]

puisque $M(S)(I)\simeq\deleted{\unsure{H}}M(I)\deleted{\unsure{\}}}(S)\unsure{=|\simeq}\underline{\mathrm{Hom}}(S^\circ,M(I))$,
